\scanpage{49}
\mspage{20}
diagramme d'inclusions \inserted{(croissantes)} d'ens tot. ordonnés
\[
(31)
\]
\diagram{\begin{tikzcd}[row sep=large,column sep=large]
&d''=d\cup d'&\\
d\arrow[ur,hook,no head]&&d'\arrow[ul,hook',no head]\\
&\delta=d\cap d'\arrow[ul,hook',no head]\arrow[ur,hook,no head]&
\end{tikzcd}}
tel que ce soit un diagramme cocart. dans la cat. des ens. ordonnés (p.ex $\delta$ segment initial de $d$, et $d'$ contenant l'élément final de $d$) alors le diagramme correspondant
\[
(32)
\]
\diagram{\begin{tikzcd}[row sep=large,column sep=large]
&\Sigma^{\unsure{<|*}}(d'')\arrow[dl]\arrow[dr]&\\
\Sigma^{\unsure{<|*}}(d)\arrow[dr]&&\Sigma^{\unsure{<|*}}(d')\arrow[dl]\\
&\Sigma^{\unsure{<|*}}(\delta)&
\end{tikzcd}}
est cartésien. \deleted{(\unsure{pour cela} le cas est}

Les conditions b c d du n\textsuperscript{o} précédent \unsure{deviennent|s'écrivent} ici en une seule (du type d)

\textcircled{b} Considérons un ens \deleted{\unsure{ens tot}} ordonné \inserted{de ``dim'' $\geq1$} $d=(i_0,-i_n)$,
\[
\deleted{\unsure{\Sigma(n)}\longrightarrow\Sigma(\unsure{0|I})}
\]
et considérons le diagramme
\[
(33)
\]
\diagram{\begin{tikzcd}[row sep=large,column sep=large]
&\Sigma^{\unsure{<|*}}(i_0,i_n)\arrow[dr,"\varphi_0"]&\\
\Sigma^{\unsure{<|*}}(d)\arrow[ur]\arrow[r]\arrow[r,no head,shift right=0.7ex,shorten >=6pt]\arrow[ddr,bend right=15]&\Sigma^{\unsure{<|*}}(i_1,i_n)\arrow[r,"\varphi_1"]&\Sigma^{\unsure{<|*}}(i_n)\\
&\substack{\cdot\\\shortmid}&\\
&\Sigma^{\unsure{<|*}}(i_{n-1},i_n)\arrow[uur,"\varphi_{n-1}"']&
\end{tikzcd}}
\deleted{de \unsure{plongements} bordants}. d'immersions, où les $\varphi_0,\varphi_1,\ \varphi_{n-1}$ sont bordantes. On suppose
\margin{Plus précisément \unsure{nous}, on veut que pour $x\in\Sigma^{\unsure{<|*}}(0)$, les $z_\alpha\in\Sigma^{\unsure{<|*}}(1)$ au dessus de $x$ (par but) sont en \uline{corr. biunivoque} avec les germes $B_\alpha$ des sous-ens. bord.-ts, images locales des $\Sigma^{\unsure{<|*}}(1)$ en les $z_\alpha$, que ces germes sont transversaux, et que les points \inserted{de $\Sigma^{\unsure{<|*}}(1)$} \deleted{\illegible{3}} \unsure{au dessus de $x$} \unsure{correspondants} \unsure{distincts} \unsure{déterminent} un \unsure{drapeau} \unsure{dans} les $z_\alpha\mathinner{\cdotp-}$}
\margin{48}
